\chapter{\textbf{INTRODUCCIÓN}}
La medición de la pobreza en Colombia enfrenta importantes desafíos derivados de la limitada
representatividad muestral de las encuestas de hogares y de la ausencia de información actualizada con
cobertura completa para todos los municipios del país. En el caso particular del Índice de Pobreza
Multidimensional (IPM), la última medición con alcance municipal proviene del Censo de 2018, lo que
genera un vacío significativo para el análisis territorial y la focalización de políticas públicas. Esta carencia dificulta identificar desigualdades persistentes, evaluar agendas como el avance en los Objetivos de Desarrollo Sostenible y orientar intervenciones sociales de manera precisa. \\

\setlength{\parindent}{2em} En este contexto, la estimación para áreas pequeñas (Small Area Estimation, SAE) se presenta como una alternativa metodológica adecuada para generar indicadores robustos en dominios con escasa o nula representatividad muestral. La posibilidad de integrar información censal, administrativa y muestral, así como de incorporar patrones de dependencia geográfica mediante modelos espaciales, permite formular estimaciones coherentes y territorialmente consistentes. Estas características hacen que la combinación de modelos SAE y enfoques espaciales sea especialmente pertinente para el caso colombiano, donde las brechas territoriales requieren instrumentos analíticos más finos y oportunos. \\

\setlength{\parindent}{2em} El proyecto se orienta a desarrollar una metodología estadística de estimación municipal del IPM que articule estas metodologías con el uso de fuentes oficiales del DANE y otras bases de datos
complementarias (también abiertas). Asimismo, se propone evaluar diferentes alternativas de modelamiento, contrastarlas a partir de criterios de coherencia geográfica y desempeño estadístico, y analizar los patrones territoriales que surjan del ejercicio. \\

\setlength{\parindent}{2em} En conjunto, esta propuesta busca contribuir a mejorar la disponibilidad de información para el diseño y focalización de políticas públicas, fortalecer el análisis territorial y aportar una herramienta metodológica replicable para el estudio de fenómenos sociales en contextos de baja representatividad estadística.